\documentclass{article}
\usepackage{amsmath,amssymb,amsthm}
\newtheorem{theorem}{Theorem}
\begin{document}
\title{Complementary-rail lifting does not supply an affine selector}
\author{}
\date{November 24, 2026}
\maketitle

\paragraph{Status.}
The obstruction below is Lean checked in
\texttt{JointFamilyDualRailPenalty.lean}. It does not modify or contradict
the inherited exact lifting theorem in \texttt{DualRailLift.lean}.
It rules out one proposed use of that lift in the affine-penalty gluing
construction, not nonlinear selection or the original upper-bound target.

\begin{theorem}
Let $L$ be an affine function nonnegative on the Boolean cube. If
$L(x)=L(1-x)=0$ for one antipodal pair, then $L$ vanishes everywhere
on the cube. Consequently, if $L$ vanishes on all complementary-rail
points $(x,1-x)$, it cannot suppress any other Boolean input.
\end{theorem}
\begin{proof}
Writing $L(y)=c+\sum_i b_i y_i$ gives
\[
 L(y)+L(1-y)=2c+\sum_i b_i=L(x)+L(1-x)=0.
\]
Both terms are nonnegative, so each is zero. Complementing every
coordinate of $(x,1-x)$ gives another complementary-rail point.
The Lean declarations are
\texttt{jointFamily\_nonnegative\_affine\_antipodal\_zeros} and
\texttt{jointFamily\_dualRail\_affine\_penalty\_zero}.
\end{proof}

Already on one pair, the affine parallelogram identity is
$L(00)+L(11)=L(01)+L(10)$. Nonnegativity and zeros at $01,10$ force
zeros at $00,11$. This is separately checked as
\texttt{jointFamily\_nonnegative\_affine\_pair\_obstruction}.

\paragraph{What the matching cover would have supplied.}
In a layer of weight $k$ on $2k$ coordinates, a perfect matching
selects the $2^k$ points having exactly one one in each pair. Each
layer point belongs to exactly $k!$ matchings, out of $(2k-1)!!$
total matchings. This regular incidence gives an explicit fractional
cover and standard random partial covers. These are combinatorial
facts only: a score fitting one matching cube still has uncontrolled
signs at other layer points. The affine penalty used by the existing
global transfer cannot remove those values, by the theorem above.

On the fixed weight-$k$ layer, the number of violated matching pairs
equals twice the number of pairs equal to $11$, so a natural penalty
is a sum of quadratic products. Implementing its conjunction or
composition by a small sum of admissible affine ratios is a separate
problem. Expanding a product over all $k$ pairs gives an exponential
cost and supplies no saving. A positive continuation must either
control signs across the entire layer, construct a cheaper nonlinear
selector, or change the gluing geometry.

\end{document}
